\section{Introduction}\label{sec:intro}

\begin{figure}[t]
\centering
\resizebox{\linewidth}{!}{%
\begin{tikzpicture}[
  font=\footnotesize,
  blk/.style={draw=#1!70!black, fill=#1!10, rounded corners=2pt, thick,
    text width=2.75cm, minimum height=1.0cm, align=center, inner sep=3pt},
  blk/.default=blue,
  io/.style={draw=gray, dashed, rounded corners=2pt, text width=1.75cm,
    minimum height=0.8cm, align=center, inner sep=2pt},
  arr/.style={-{Stealth[length=5pt]}, thick, draw=black!75},
  sig/.style={-{Stealth[length=5pt]}, thick, dashed, draw=violet!80!black},
  lab/.style={font=\scriptsize, fill=white, inner sep=1pt},
]
% inputs
\node[io] (plan) at (0, 1.45) {global plan\\$\sigma$};
\node[io] (odom) at (0, 0) {odometry\\$\mathbf{x}_t$};
\node[io] (cmap) at (0,-1.45) {local\\costmap};
% controller internals
\node[blk=violet] (crit) at (3.35, 1.45) {\textbf{EscapeCritic}\\APF + gap attraction,\\injected only if $e_t{=}1$};
\node[blk] (mppi) at (7.05, 1.45) {\textbf{MPPI optimizer}\\(Nav2, reused)\\$K$ rollouts, softmax};
\node[blk=violet] (det) at (3.35, 0) {\textbf{Entrapment detector}\\stall of $k^{\max}_t$ $\Rightarrow e_t$};
\node[blk=orange] (trk) at (3.35,-1.45) {\textbf{Obstacle tracker}\\clusters $\to(\mathbf{p}_o,R_o,\mathbf{v}_o)$};
\node[blk=violet] (crd) at (7.05,-0.72) {\textbf{Coordinator}\\$\alpha_t$ from $e_t$ and TTC};
\node[blk=teal] (cbf) at (10.75, 0) {\textbf{CBF-QP filter}\\look-ahead point,\\OSQP, brake if $\delta{>}\epsilon$};
\node[io, text width=2.0cm] (cm) at (13.9, 0) {\texttt{collision\_}\\\texttt{monitor} $\to$ base};
% container
\begin{scope}[on background layer]
\node[draw=black!60, fill=black!2, rounded corners=4pt, inner sep=7pt,
  fit=(crit)(mppi)(det)(trk)(crd)(cbf),
  label={[font=\footnotesize\bfseries, anchor=south west, inner sep=2pt]north west:\sname\ \texttt{SafeEscapeController} (\texttt{nav2\_core::Controller} plugin)}] (box) {};
\end{scope}
% edges
\draw[arr] (plan) -- (crit);
\draw[arr] (plan.north) -- ++(0,0.55) -| node[lab, pos=0.25] {path, goal} (mppi.north);
\draw[arr] (odom) -- node[lab, above] {pose} (det);
\draw[arr] (cmap) -- (trk);
\draw[arr] (crit) -- node[lab, above] {costs} (mppi);
\draw[sig] (det) -- node[lab, right] {$e_t$} (crit);
\draw[sig] (det.east) -- ++(0.35,0) |- node[lab, pos=0.75, above] {$e_t$} (crd.west);
\draw[arr] (trk.east) -- ++(0.35,0) |- (crd.west |- 0,-1.0);
\draw[arr] (trk.south) -- ++(0,-0.35) -| node[lab, pos=0.25] {dynamic obstacles only} (cbf.south);
\draw[arr] (mppi.east) -| node[lab, pos=0.3, above] {$\mathbf{u}^{\mathrm{nom}}_t$} (cbf.north);
\draw[arr] (crd.east) -- ++(0.3,0) |- node[lab, pos=0.7, below] {$\alpha_t$} (cbf.west);
\draw[arr] (cbf) -- node[lab, above] {$\mathbf{u}^{\mathrm{safe}}_t$} (cm);
\end{tikzpicture}}
\caption{\textbf{Overview of \sname.} The stock Nav2 MPPI optimizer produces a nominal command from sampled rollouts. A single entrapment signal $e_t$, computed from monotone progress along the global path (\cref{sec:detect}), switches the \texttt{EscapeCritic}'s repulsion and gap-attraction costs into the rollout scoring (\cref{sec:escape}) and tells the coordinator to raise the CBF gain $\alpha_t$ unless a tracked obstacle's time-to-collision is imminent (\cref{sec:coord}). The CBF-QP filter then projects the nominal command onto the barrier-safe set for the tracked dynamic obstacles (\cref{sec:cbf}); Nav2's \texttt{collision\_monitor} remains the last reactive layer.}
\label{fig:overview}
\end{figure}


Autonomous mobile robots increasingly work in cluttered, dynamic indoor spaces such as warehouses, hospitals and service floors, where the local controller must be fast, non-conservative and safe at the same time. Within the ROS~2 Nav2 ecosystem \citep{nav2}, the MPPI controller (\code{nav2_mppi_controller}) is a widely used high-performance option: it samples a batch of control rollouts, scores them with a bank of critic plugins, and fuses them with an information-theoretic softmax \citep{mppi_orig}. With the defaults we build on, it plans over a horizon of $56\times0.05\,\mathrm{s}\approx2.8\,\mathrm{s}$ with $1{,}000$ rollouts per cycle.

Two weaknesses recur in exactly the environments that matter most. First, \emph{local-minima entrapment}: a finite horizon and zero-mean Gaussian sampling make the robot stall or oscillate in front of non-convex obstacles. Within the MPPI critic set, Nav2's mitigations (\code{PreferForwardCritic}, \code{TwirlingCritic}) are always-on cost shaping that neither detects entrapment nor steers an escape. At the stack level, Nav2 does detect a stall: the controller server's progress checker reports a failure when the robot has not moved a set distance within a set time, and the default behaviour tree then clears the costmaps and runs generic recoveries such as spinning, waiting and backing up \citep{nav2,nav2_progress,nav2_bt}. These recoveries are not directed toward an opening. In the MPPI literature, escape methods either reshape the sampling distribution all the time \citep{log_mppi,tsallis_mppi,svg_mppi,biased_mppi,flow_mppi} or add guidance: GP-MPPI recommends a free-space subgoal from a learned navigability surface \citep{gp_mppi}, and repulsive-potential augmentation adds an artificial potential field to the MPPI cost \citep{drpa_sii}, which DRPA-MPPI switches on only when it detects entrapment \citep{drpa_mppi}. None of these carries a safety certificate.

Second, \emph{no formal safety}: in Nav2 MPPI, collision avoidance is encoded as soft costs, the costmap is a present-time snapshot with no prediction of obstacle motion, and \code{collision_monitor} is purely reactive. CBF--MPPI fusions supply safety \citep{shield_mppi,br_mppi,scbf_mppi,dualguard,gs_mppi}, but none of them detects entrapment or adds an escape mechanism (GS-MPPI counters the myopic behaviour often associated with CBFs through long-term performance terms), and none ships as a Nav2 controller. Among CBF integrations with Nav2 that we could verify, the Nav2 documentation describes inserting a third-party product, the 3Laws Supervisor, which its vendor describes as based on CBFs, into the command chain after the controller \citep{nav2_3laws,threelaws_pr}; it is a separate node rather than a controller plugin, and neither source describes a mechanism for escaping local minima.

The two weaknesses also interact: conservatism for safety invites entrapment, and aggressive escape invites collision. Several works outside Nav2 and MPPI address both at once, among them model predictive control that escapes local minima while keeping a safe backup trajectory \citep{soloperto}, gap-based local planners with provable collision-free properties \citep{potential_gap,safer_gap,dynamic_gap}, and modulated CBF filters \citep{mcbf_qp} that reduce or eliminate, among moving obstacles, the spurious equilibria that CBF quadratic programs (QPs) can introduce \citep{reis_cbf_equilibria}, with concurrent extensions \citep{xue_proactive,chen_milp_cbf}.

We present \emph{\sname}, a Nav2 controller plugin that subclasses the stock MPPI controller and post-processes its command each cycle (\cref{fig:overview}). A single entrapment detector, driven by monotone progress along the global path and running inside the controller rather than at the stack level, is the one source of truth for the whole controller. When it fires, an \code{EscapeCritic} switches repulsion and free-space gap-attraction costs into the rollout scoring, and a coordinator raises the gain of a look-ahead-point CBF filter so that the escape manoeuvre is admitted, unless a tracked obstacle's time-to-collision is imminent. This is a system contribution, and we state it as such:
\begin{itemize}
    \item \emph{A Nav2 Controller Integrating Escape and CBF Safety}: a public Nav2 local-controller plugin that combines online local-minima detection-and-escape with a CBF safety filter for tracked dynamic obstacles, packaged as a \code{nav2_core::Controller} plus a pluginlib MPPI critic over a plugin-free core, and released under Apache-2.0. Each ingredient has precedent: gated repulsion in DRPA-MPPI \citep{drpa_mppi}, free-space subgoal and gap guidance in GP-MPPI \citep{gp_mppi} and gap-based planners \citep{potential_gap,safer_gap}, CBF filtering of MPPI commands \citep{scbf_mppi}, and a supervisor for Nav2 that its vendor describes as CBF-based \citep{nav2_3laws,threelaws_pr}. What is new is their integration inside one Nav2 controller; to our knowledge no public controller plugin already provides it.
    \item \emph{An Entrapment-Gated Gain Schedule}: a rule that raises the CBF class-$\mathcal{K}$ gain on detected entrapment and drops it back on an imminent time-to-collision. Online adaptation of CBF gains is established, including selecting class-$\mathcal{K}$ parameters to reduce predicted deadlock \citep{kim_adaptive_cbf} and letting MPPI optimize a barrier-rate parameter \citep{br_mppi}; our rule is a simpler, rule-based two-level instance of online gain adaptation whose specific elements are its trigger and its time-to-collision override. Its invariance property (\cref{prop:invariance}) is a known result for time-varying class-$\mathcal{K}$ terms \citep{ada_cbf,rt_cbf}, which we restate as a conditional lemma: it covers the look-ahead point and the tracked dynamic obstacles, assumes the barrier inequality holds between samples, and requires a slack-free QP. In our benchmark the schedule is empirically indistinguishable from an independent escape+CBF stack.
    \item \emph{Paired Benchmark and Deployment Evidence}: a randomized, seed-deterministic $1{,}200$-trial benchmark on a Python 2D mirror of the controller's primitives ($200$ seeded worlds, each run under six configurations) with paired McNemar tests and Holm correction, every number regenerated from committed per-trial data, together with a live ROS~2 Jazzy, Nav2 and Gazebo integration that exposed three faults unreachable from unit tests.
\end{itemize}

We are deliberately conservative about novelty and about what we claim empirically. Fusing CBFs with MPPI is not new, and neither is gated escape, online adaptation of the CBF gain, or treating escape and safety together; the paper is best read as a system report on putting these pieces into one Nav2 controller. The measured win, in the 2D mirror, is the online detection-and-escape layer: success rises from $0\%$ to $88$--$90\%$ on the static U-trap family and from $0\%$ to $62$--$78\%$ on the narrow-dynamic family, and the gap search is load-bearing, since without it both families stay at $0\%$ (\cref{sec:results,sec:ablation}). In the mirror the gap acts as a temporary subgoal, close to the subgoal guidance of GP-MPPI \citep{gp_mppi}; the gap cost of the C++ critic has not been ablated. The coordination has the conditional invariance property of \cref{prop:invariance} yet produces no measurable outcome difference in the regimes we could test, and we delimit where it should matter (\cref{sec:discussion}). The remaining claims are Nav2-native deployment and the faults a live stack exposes; no live run has yet reached its goal, and no stack-level or third-party baseline has been run (\cref{sec:discussion}).
